\section{Sample Collection \& Bioanalytical Methods}
\label{sec:sampling}

\subsection{Pharmacokinetic Blood Sampling Schedule}

Serial blood samples for PK determination will be collected via indwelling intravenous cannula or direct venipuncture. A total of 16 PK blood samples will be collected from each subject over a 72-hour post-dose period. 

Table~\ref{tab:pk_schedule} outlines the exact timing and tolerance windows for each blood collection point relative to the time of oral dosing ($T_0 = 0.0$ hours).

\begin{table}[H]
\centering
\small
\caption{PK Blood Collection Schedule and Sampling Windows}
\label{tab:pk_schedule}
\begin{tabularx}{\textwidth}{@{}l l l X@{}}
\toprule
\textbf{Time Point} & \textbf{Target Time (h)} & \textbf{Window (min)} & \textbf{Rationale} \\
\midrule
Pre-dose & -0.50 (within 30m prior) & -- & Establish baseline level \\
Post-dose & 0.25 & $\pm$ 2 & Capture initial absorption phase \\
Post-dose & 0.50 & $\pm$ 2 & Capture absorption phase \\
Post-dose & 1.00 & $\pm$ 5 & Capture near-peak levels \\
Post-dose & 1.50 & $\pm$ 5 & Capture peak concentration range \\
Post-dose & 2.00 & $\pm$ 5 & Expected $t_{\max}$ range \\
Post-dose & 3.00 & $\pm$ 10 & Capture early distribution phase \\
Post-dose & 4.00 & $\pm$ 10 & Capture distribution phase \\
Post-dose & 6.00 & $\pm$ 10 & Capture early elimination phase \\
Post-dose & 8.00 & $\pm$ 15 & Capture elimination phase \\
Post-dose & 12.00 & $\pm$ 15 & Capture elimination phase \\
Post-dose & 24.00 & $\pm$ 30 & Confirm 24h trough concentration \\
Post-dose & 36.00 & $\pm$ 45 & Extended elimination phase \\
Post-dose & 48.00 & $\pm$ 60 & Extended elimination phase \\
Post-dose & 72.00 & $\pm$ 120 & Determine terminal elimination phase \\
\bottomrule
\end{tabularx}
\end{table}

\subsection{Urine Sampling Schedule}

Urine samples will be collected cumulatively over the following intervals post-dose: pre-dose (-2 to 0 hours), 0--4 hours, 4--8 hours, 8--12 hours, 12--24 hours, and 24--48 hours. Total urine volume will be recorded for each interval, and a 10~mL aliquot will be retained and frozen at $-80^{\circ}$C for assaying parent drug and metabolite concentrations.

\subsection{Bioanalytical Methodology}

Plasma and urine concentrations of SYN-409 will be determined using a validated high-performance liquid chromatography-tandem mass spectrometry (LC-MS/MS) method \citep{zhao2025validation}. 

\subsubsection{Sample Processing}
Blood samples will be collected in 4~mL K$_2$EDTA tubes and immediately inverted 8--10 times. Tubes will be centrifuged at 1500~g for 10 minutes at $4^{\circ}$C within 30 minutes of collection. The resulting plasma will be split into two aliquots (primary and back-up) and stored in polypropylene tubes at $-80^{\circ}$C.

\subsubsection{LC-MS/MS Assay Specifications}
\begin{itemize}
  \item \textbf{Calibration Range:} 0.500 to 1000~ng/mL in plasma; 2.00 to 5000~ng/mL in urine.
  \item \textbf{Lower Limit of Quantification (LLOQ):} 0.500~ng/mL.
  \item \textbf{Internal Standard:} Deuterated SYN-409 ($d_4$-SYN-409).
  \item \textbf{Precision and Accuracy:} Within-run and between-run coefficient of variation (CV) $< 15\%$ ($< 20\%$ at LLOQ), and accuracy within $85\%$ to $115\%$ of nominal values.
\end{itemize}
